% New author continuation of native330--332. No external spectral theorem is invoked.
\section{Absolute Fourier inversion of the unchanged microscopic intervals}
\label{sec:ui-exact}

The weak roof estimate in Section~\ref{sec:wo-endpoint-tests} has a
nonintegrable $|s|^{-1}$ tail when it is used alone.  The transform of an
\emph{exact bounded roof interval} supplies another inverse power.  We
exploit this factor without replacing the interval by a likelihood or
enlarging it.  The resulting formula is an absolutely convergent inversion
of the actual event.  Its error after division by the true interval
length is stated explicitly.

The earlier exact reconstruction arguments are retained. Historical
native 64 already gives a sharp full-law interval integral with an
extracted second-derivative budget; native 67 gives absolute sharp
inversion for the stationary flight-age overlap; native 93 gives
coefficient-first raw reconstruction after its exact germ subtraction.
Theorem~\ref{thm:av-likelihood} already permits different windows
at different labels, and Corollary~\ref{cor:av-denominator} already
uses their actual selected denominator for positive likelihoods.
No priority is claimed for Fourier inversion, interval transforms,
label summation or denominator certification as general methods.
The addition here is a direct full-joint interval corollary of the
complete first-order BV budget, including its precise sharp tail,
the original $m^2/|I|$ payment and specified endpoint tests, with
neither a remaining likelihood kernel nor a germ subtraction.

Fix $R\in[R_-,R_+]$, the actual collision count $m\ge1$, and either
$w=1$ or the original $w=W^\varepsilon$ with $\varepsilon>0$ fixed.
Use the endpoint tests and physical orbit phase of
\eqref{eq:wo-measures}, and abbreviate
\[
 V(x)=a(x)b(T_R^mx)e^{i\Psi_m(x)},\qquad
 A_V=a_0b_0,\qquad D_V=\mathcal V_m^w(a,b,\Psi).
\]
The phase is optional.  All tests have precisely the $C^1$ norms stated
there.  Write $\mathbf j=(k_1,k_2,n)\in\mathbb Z^3$ and
$f_{\mathbf j}=f^{w;a,b,\Psi}_{(n,k),m,R}$; this is only an ordering
of the unchanged discrete record.  Define its genuine joint transform
\begin{equation}\label{eq:ui-transform}
 \begin{split}
 \mathcal F_{m,R}^{w,V}(\vartheta,s)
 &=\sum_{\mathbf j}e^{i\vartheta\cdot\mathbf j}
                         \int_{\mathbb R}e^{ist}f_{\mathbf j}(t)\,dt\\
 &=\int w\eta_R(T_R^mx)V(x)
       e^{i\{\xi\cdot K^{\rm c}_{m,R}(x)+\zeta A_{m,R}(x)
                         +sL_{m,R}(x)\}}\,d\nu_R^*(x).
 \end{split}
\end{equation}
Here $\vartheta=(\xi_1,\xi_2,\zeta)\in[-\pi,\pi]^3$.
Theorem~\ref{thm:wo-bv} gives
\begin{equation}\label{eq:ui-input}
 \sum_{\mathbf j}\|f_{\mathbf j}\|_1\le A_V,\qquad
 \sum_{\mathbf j}\|Df_{\mathbf j}\|_{\rm TV}\le D_V,
 \qquad
 |\mathcal F_{m,R}^{w,V}(\vartheta,s)|
       \le\min\{A_V,D_V/|s|\}\quad(s\ne0).
\end{equation}
At fixed $m$ only finitely many discrete records are attainable, by the
physical finite horizon.  None of the following constants uses their
number.

Let $I$ be a union of at most $J\ge1$ bounded intervals, let
$d=|I|$, and use the convention
$\widehat{\mathbf1_I}(s)=\int_Ie^{ist}\,dt$.  Endpoints may be open
or closed, because the true roof laws in question have densities.
Merging overlapping intervals gives disjoint intervals without
increasing $J$, and direct integration gives
\begin{equation}\label{eq:ui-window-transform}
 |\widehat{\mathbf1_I}(s)|\le\min\{d,2J/|s|\}\quad(s\ne0).
\end{equation}
For $S>0$ define
\begin{equation}\label{eq:ui-cutoff}
 \mathfrak I_S(\mathbf j,I;V)
  =(2\pi)^{-4}\int_{[-\pi,\pi]^3}\int_{-S}^{S}
      e^{-i\vartheta\cdot\mathbf j}\widehat{\mathbf1_I}(-s)
           \mathcal F_{m,R}^{w,V}(\vartheta,s)\,ds\,d\vartheta.
\end{equation}
This is a sharp frequency truncation, and is not asserted positive.

\begin{theorem}[Exact interval inversion and its full tail]
\label{thm:ui-inversion}
For every specified physical parameter, exact label and interval union,
\begin{equation}\label{eq:ui-inversion}
 \int_I f_{\mathbf j}(t)\,dt
  =(2\pi)^{-4}\int_{[-\pi,\pi]^3}\int_{\mathbb R}
      e^{-i\vartheta\cdot\mathbf j}\widehat{\mathbf1_I}(-s)
            \mathcal F_{m,R}^{w,V}(\vartheta,s)\,ds\,d\vartheta.
\end{equation}
The integral on the right is absolutely convergent.  If $d>0$, then
\begin{align}
 \left|\int_I f_{\mathbf j}-\mathfrak I_S(\mathbf j,I;V)\right|
 &\le {D_V\over\pi}\,\mathfrak t_J(d,S)
 \le {2JD_V\over\pi S},                         \label{eq:ui-tail}\\
 \mathfrak t_J(d,S)
 &=\begin{cases}
 d\bigl[1+\log\bigl(2J/(dS)\bigr)\bigr],&dS\le2J,\\
 2J/S,&dS\ge2J.
 \end{cases}                                    \label{eq:ui-tail-function}
\end{align}
For $d=0$ the event and every displayed integral vanish.  The assertions
hold for all original lattice and occupation labels, including zero
arithmetic classes, and for the full physical roof support.  There is
exactly one conditioning factor $c_*^{-1}$, already present in
$f_{\mathbf j}$ and $D_V$.
\end{theorem}
\begin{proof}
For bounded $s$ the integrand in \eqref{eq:ui-inversion} is bounded by
$dA_V$.  At $|s|\ge1$, its absolute value is bounded by
$2JD_V/s^2$.  The torus has finite volume; these two estimates prove
absolute integrability.

For completeness we justify inversion without assuming pointwise
Fourier inversion for a BV density.  Let $k_{\mathfrak b}$ be the
probability kernel of Lemma~\ref{lem:av-kernel}.  The BV translation
bound implies
\[
 \|k_{\mathfrak b}*f_{\mathbf j}-f_{\mathbf j}\|_1
   \le\|Df_{\mathbf j}\|_{\rm TV}
                  \kappa_{\rm ker}/\mathfrak b\longrightarrow0.
\]
For this intermediate convolution, the compact-frequency formula and
orthogonality on the entire three-dimensional torus give the right
side of \eqref{eq:ui-inversion} with the additional factor
$\eta(s/\mathfrak b)$.  Every integral is then over a bounded
frequency domain, so Fubini is valid.  The left side tends to
$\int_I f_{\mathbf j}$ by the displayed $L^1$ estimate.
Since $|\eta|\le1$ and $\eta(s/\mathfrak b)\to1$, the already
proved integrable bound permits dominated convergence on the right.
The convolution has disappeared from the resulting identity.  Its
left side is the exact original interval event, with its original
width and trajectory information.

The three torus integrals cancel three of the four factors $2\pi$
in an absolute tail bound.  Equations~\eqref{eq:ui-input} and
\eqref{eq:ui-window-transform} therefore yield
\[
 \left|\int_I f_{\mathbf j}-\mathfrak I_S(\mathbf j,I;V)\right|
 \le {D_V\over\pi}\int_S^\infty
                       \min\{d/u,2J/u^2\}\,du.
\]
Splitting this positive integral at $2J/d$ proves
\eqref{eq:ui-tail-function}.  Its integrand is always at most
$2J/u^2$, proving the last bound in \eqref{eq:ui-tail} for every
$S>0$, including $dS<2J$.  No roof supremum, label supremum,
collision limit or outer geometric limit was exchanged.
\end{proof}

\begin{corollary}[All labels and their different exact windows]
\label{cor:ui-all-labels}
For each attainable $\mathbf j$ choose an empty set or a union
$I_{\mathbf j}$ of at most $J$ bounded intervals, and choose
coefficients $|c_{\mathbf j}|\le1$.  Summing the exact component
identities gives
\begin{equation}\label{eq:ui-all-tail}
 \left|\sum_{\mathbf j}c_{\mathbf j}
       \left(\int_{I_{\mathbf j}}f_{\mathbf j}
           -\mathfrak I_S(\mathbf j,I_{\mathbf j};V)\right)\right|
 \le {2JD_V\over\pi S}.
\end{equation}
The bound has no factor counting records, selected windows or
components of their physical preimages.
\end{corollary}
\begin{proof}
Integrate the full torus first in \eqref{eq:ui-cutoff}; orthogonality
selects exactly $\widehat f_{\mathbf j}(s)$.  For that component the
same proof of \eqref{eq:ui-tail} gives
$2J\|Df_{\mathbf j}\|_{\rm TV}/(\pi S)$.  Sum these positive
error bounds and apply \eqref{eq:ui-input}.  A finite count has
finitely many attainable labels, so no infinite complex-series
interchange is needed.  The norm sum explains why the estimate is
independent of that finite number.
\end{proof}

\begin{corollary}[The original density normalization]
\label{cor:ui-normalization}
For a nonempty original interval union $I$, of its actual length $d$,
\begin{equation}\label{eq:ui-normalized-tail}
 {m^2\over d}
 \left|\int_I f_{\mathbf j}-\mathfrak I_S(\mathbf j,I;V)\right|
 \le {2JD_Vm^2\over\pi dS}.
\end{equation}
In particular, for any $P>0$, the analytical choice
$S=m^{P+2}/d$ makes this upper bound
$2JD_Vm^{-P}/\pi$.  This holds for every positive $d$, however
small.  A claim at this bandwidth must evaluate the joint transform
on that whole bandwidth; no available low-frequency theorem is
silently extended to it.
\end{corollary}
\begin{proof}
Divide \eqref{eq:ui-tail} by the true interval length and multiply by
the original $m^2$.  Substitute the stated finite $S$.  There is no
lower bound on $d$ and no modification of the event.
\end{proof}

\subsection{Conditional observations use the actual event mass}

Let $\mathcal A$ be one exact record with its interval union, or an
event formed by the different exact windows in
Corollary~\ref{cor:ui-all-labels}, with coefficients all equal to one.
Use the positive physical source $w\eta_R(T_R^mx)d\nu_R^*$ and
write
\[
 p=\int\mathbf1_{\mathcal A}w\eta_R(T_R^mx)d\nu_R^*,\qquad
 N=\int V\mathbf1_{\mathcal A}w\eta_R(T_R^mx)d\nu_R^*.
\]
Let $D_S$ and $N_S$ be the corresponding sharp-frequency expressions
with test $1$ and $V$, respectively.  The former is real by conjugate
symmetry.  Put $E_0=2JD_{\rm all}^w/(\pi S)$ and
$E_V=2JD_V/(\pi S)$; the smaller single-window bounds in
\eqref{eq:ui-tail} may be substituted when appropriate.

\begin{proposition}[Certified denominator and conditional test]
\label{prop:ui-conditional}
One has $|p-D_S|\le E_0$ and $|N-N_S|\le E_V$.
If $D_S>E_0$, then the actual selected probability satisfies
$p\ge D_S-E_0>0$, and
\begin{equation}\label{eq:ui-conditional}
 \left|{N\over p}-{N_S\over D_S}\right|
       \le {E_V+A_VE_0\over D_S}.
\end{equation}
No conditional law is asserted when the actual event has zero mass.
This statement concerns exactly the displayed smooth endpoint and
trajectory-phase observations.  It does not assign their variation
bound to an arbitrary additional Borel information selector.
\end{proposition}
\begin{proof}
The two absolute errors follow from Theorem~\ref{thm:ui-inversion}
and Corollary~\ref{cor:ui-all-labels}.  The positivity certificate
follows from the first error.  Since $|V|\le A_V$ on the original
trajectories, positivity gives $|N|\le A_Vp$.  The exact identity
\[
 {N\over p}-{N_S\over D_S}
       ={N-N_S\over D_S}+{N(D_S-p)\over pD_S}
\]
proves \eqref{eq:ui-conditional}.  This uses the mass of precisely
$\mathcal A$, not that of a containing event.  In particular,
if $p=0$ the positivity certificate cannot hold because
$|D_S|\le E_0$.
\end{proof}

\subsection{A sufficient pointwise bridge, with the missing input exposed}

We state a sufficient implication to identify what the new interval
formula does and does not complete.  It retains the existing raw
arithmetic reference, all its peaks, and the ordered outer limit.
Here $w=1$, $V=1$, and $p_{\mathbf j,m,R}$ is the original density.
Every raw-target supremum below uses the original ranges
$R\in[R_-,R_+]$, $1\le n\le m$ and $k\in\mathbb Z^2$.
Set
\[
 g_{\mathbf j,m,R}(u)=m^{-2}\mathcal L_{m,R}(k_1,k_2,u,n).
\]
Fix $M$.  At each outer stage $B$ choose a positive $h_B$ before
the collision limit, and for every central $t$ put
$I_{B,t}=(t-h_B,t+h_B)$.  These intervals are analysis tests of the
same physical roof, not a redefinition of any conditioning event.
For any finite bandwidth $S_{m,B}>0$ define
\begin{align*}
 \Omega_{m,B}^{p}
   &=\sup_{R,n,k}m^2\esssup_{t\in\mathcal C_{m,R,n,k}(M)}
                       |Dp_{\mathbf j,m,R}|([t-h_B,t+h_B]),\\
 \Omega_{m,B}^{g}
   &=\sup_{R,n,k}\esssup_{t\in\mathcal C_{m,R,n,k}(M)}
      {1\over2h_B}\int_{t-h_B}^{t+h_B}
       |\mathcal L_{m,R}(k_1,k_2,u,n)
                  -\mathcal L_{m,R}(k_1,k_2,t,n)|\,du,\\
 \Xi_{m,B}
   &=\sup_{R,n,k}\esssup_{t\in\mathcal C_{m,R,n,k}(M)}
      {m^2\over2h_B}
      \left|\mathfrak I_{S_{m,B}}(\mathbf j,I_{B,t};1)
                          -\int_{I_{B,t}}g_{\mathbf j,m,R}(u)\,du\right|.
\end{align*}
The derivative measure is that of the complete actual exact-label
density, after every physical word and boundary has been included.

\begin{proposition}[An ordered raw-density implication]
\label{prop:ui-raw-bridge}
At every finite $m,B$,
\begin{equation}\label{eq:ui-raw-bridge}
 \sup_{R,n,k}\esssup_{t\in\mathcal C_{m,R,n,k}(M)}
     |m^2p_{\mathbf j,m,R}(t)-\mathcal L_{m,R}(k_1,k_2,t,n)|
 \le\Omega_{m,B}^{p}+\Omega_{m,B}^{g}+\Xi_{m,B}
               +{D_{\rm all}^{1}m^2\over\pi h_BS_{m,B}}.
\end{equation}
Consequently the original raw target follows if each of the four
terms on the right tends to zero under
$\lim_{B\to\infty}\limsup_{m\to\infty}$.  With
$S_{m,B}=m^{P+2}/h_B$, $P>0$, the last term already tends to zero
for every fixed $B$.  The local derivative-measure condition,
the reference modulus and the scalar band comparison are displayed
separately in this general inequality. The following subsection
uses the retained fixed-window theorem to pay the latter two;
they must not be relisted as new unproved requirements when that
inherited theorem is invoked at its original scope.
\end{proposition}
\begin{proof}
For a BV representative at a Lebesgue point $t$, its difference
from almost every value on $I_{B,t}$ is bounded by
$|Dp|([t-h_B,t+h_B])$.  Integrating gives
\[
 \left|p(t)-{1\over2h_B}\int_{I_{B,t}}p(u)\,du\right|
       \le |Dp|([t-h_B,t+h_B]).
\]
Insert the exact interval average between the raw density and the
reference, then insert its sharp-frequency expression and the
reference interval average.  These four differences are bounded,
respectively, by $\Omega^p$, the normalized tail
\eqref{eq:ui-normalized-tail} with $J=1,d=2h_B$, $\Xi$, and
$\Omega^g$.  This proves \eqref{eq:ui-raw-bridge}, first for each
fixed radius and almost every roof, and then with the stated suprema.
Take the collision limsup with $h_B$ already fixed; take the outer
limit last.  No interchange is needed.  The displayed bandwidth
reduces its paid term to $D_{\rm all}^1m^{-P}/\pi$.
\end{proof}

\subsection{Evaluation of the bridge by the retained fixed-window theorem}

The preceding sufficient inequality must be read together with the
written fixed-window theorem already present in historical native 87.
Its Theorem~\textup{v41-uniform-transition}, equation
\textup{v41-uniform-interval}, concerns precisely the event
\eqref{eq:v76-event} and the finite-peak reference
\eqref{eq:v76-reference}. It includes the moving nonunit branches;
it is not restricted to one fixed radius or to a unit arithmetic
factor. We retain that theorem with its existing analytic hypotheses.
The following implication uses its conclusion, without reasserting an
independent reconstruction of its whole spectral proof.

For a fixed $h>0$ set
\begin{equation}\label{eq:ui-inherited-window-error}
 \begin{aligned}
 \mathscr E_m(h)={}&
 \sup_{\substack{R,\,1\le n\le m,\,k\in\mathbb Z^2,\,t\in\mathbb R}}
 \left|\frac{m^2}{2h}\int_{t-h}^{t+h}p_{n,R}(k,m,u)\,du\right.\\
 &\hspace{44mm}\left.
               -\mathcal L_{m,R}(k_1,k_2,t,n)\right|.
 \end{aligned}
\end{equation}
The retained fixed-window conclusion is
$\mathscr E_m(h)\to0$ for every fixed $h>0$. Its proof fixes
the spectral band before the collision limit, then removes that
band by ordered interval envelopes. It supplies no rate uniform
over collision-dependent choices of $h$.

\begin{lemma}[Uniform roof modulus of the retained reference]
\label{lem:ui-reference-modulus}
Retain the finite peak data and their uniform properties in
\eqref{eq:v76-reference} and historical native 87. Let $N_{\rm pk}$
be the number of peak charts, $B_{\rm pk}$ a common upper bound for
$|B_j(R)|$, and $a>0$ the common lower bound for
$\operatorname{Re}H_j$. For each chart use its compact parameter
support $K_j$, on a neighborhood of which its scalar data are
defined. Put
\[
 b_j(R)=(a_j(R))_3,\qquad
 \kappa_j(R)=-\log|z_j(R)|\ge0,\qquad
 \omega_m=\max_j\sup_{R\in K_j}
                 |b_j(R)|e^{-m\kappa_j(R)},
\]
where an empty supremum is zero. Then $\omega_m\to0$. With
\[
 M_0=(2\pi)^{-2}a^{-2},\qquad
 M_1=M_0\sqrt{2/(\pi a)},\qquad
 D_m=\frac{N_{\rm pk}B_{\rm pk}}{c_*}
                   \{M_0\omega_m+M_1m^{-1/2}\},
\]
one has $D_m\to0$ and
\begin{equation}\label{eq:ui-reference-derivative}
 \sup_{R,n,k,t}
 \left|\partial_t\mathcal L_{m,R}(k_1,k_2,t,n)\right|
 \le D_m.
\end{equation}
In particular $\Omega^g_{m,B}\le h_BD_m/2$ for every finite
outer stage $B$.
\end{lemma}
\begin{proof}
The retained peak theorem gives continuity on each $K_j$, bounded
roof frequencies, and $\kappa_j(R)=0\Rightarrow b_j(R)=0$.
The last implication is the original no-roof-resonance assertion;
it does not say that every continued peak has roof frequency zero.
Let $b_*<\infty$ bound all $|b_j|$. For any $\epsilon>0$, the
finite union of compact sets $\{|b_j|\ge\epsilon\}\cap K_j$
either is empty or has a strictly positive minimum
$c_\epsilon$ of $\kappa_j$. Consequently
\[
 \omega_m\le\epsilon+b_*e^{-mc_\epsilon}
\]
in the nonempty case, and $\omega_m\le\epsilon$ otherwise.
First let $m\to\infty$, then $\epsilon\downarrow0$. No rate
for $c_\epsilon$ is required.

The retained peaks are interior real maxima of
$\log|\lambda_j|$, so
$\operatorname{Re}\nabla\log\lambda_j(R,a_j)=0$.
Consequently $\mu_j=i^{-1}\nabla\log\lambda_j(R,a_j)$ is real,
as asserted in historical native 87. All Gaussian arguments below
are therefore real. For every real $z$ and every displayed complex symmetric $H$,
the defining Gaussian integral gives
\[
 |G_H(z)|\le (2\pi)^{-4}
       \int_{\mathbb R^4}e^{-a|v|^2/2}\,dv=M_0,
 \qquad
 |\partial_{z_3}G_H(z)|\le (2\pi)^{-4}
       \int_{\mathbb R^4}|v_3|e^{-a|v|^2/2}\,dv=M_1.
\]
Differentiation is justified by the same integrable Gaussian
majorant with the factor $|v_3|$. In each term of
\eqref{eq:v76-reference}, differentiating the phase costs
$|b_j|M_0$, while differentiating the Gaussian argument costs
$m^{-1/2}M_1$. Its coefficient has modulus at most
$c_*^{-1}B_{\rm pk}e^{-m\kappa_j}$. Sum the finitely many terms
and take the real part to prove
\eqref{eq:ui-reference-derivative}. This estimate is global in
the real roof target, so averaging it over the analysis interval
does not require that every point of that interval be central.
Finally
$(2h_B)^{-1}\int_{-h_B}^{h_B}|u|\,du=h_B/2$ proves the last
assertion.
\end{proof}

\begin{proposition}[The fixed-window input evaluates the scalar bridge]
\label{prop:ui-inherited-bridge}
Assume the retained fixed-window conclusion
\eqref{eq:ui-inherited-window-error} and its finite peak data at
their original scope. At each finite outer stage choose $h_B>0$
before the collision limit. For the quantities in
Proposition~\ref{prop:ui-raw-bridge}, one has
\begin{equation}\label{eq:ui-xi-inherited}
 \Xi_{m,B}\le\mathscr E_m(h_B)+\frac{h_BD_m}{2}
                +\frac{D_{\rm all}^{1}m^2}{\pi h_BS_{m,B}}.
\end{equation}
Thus, with $S_{m,B}=m^{P+2}/h_B$ and fixed $P>0$, both
$\Xi_{m,B}$ and $\Omega^g_{m,B}$ tend to zero for each fixed
$B$. More directly, writing $\mathcal E_M$ for the original
raw-error collision limsup,
\begin{equation}\label{eq:ui-local-charge-suffices}
 \mathcal E_M\le\limsup_{m\to\infty}\Omega^p_{m,B}
                         \qquad\text{for every fixed }B.
\end{equation}
In particular the sufficient remaining condition on this route is
\[
 \lim_{B\to\infty}\limsup_{m\to\infty}\Omega^p_{m,B}=0.
\]
This is a condition on the complete normalized local derivative
charge of the original physical density. It is not proved by the
global unscaled BV estimate and is not claimed necessary.
\end{proposition}
\begin{proof}
The label density here is exactly $p_{n,R}(k,m,\cdot)$, and its
interval mass is the original $n$th-return event with collision
count $m$, by \eqref{eq:v76-event}. Apply the retained fixed-window
theorem with $J=(-h_B,h_B)$. The interval length is positive and
fixed at this stage, so division by $2h_B$ is legitimate and
$\mathscr E_m(h_B)\to0$.

Insert the actual interval mass between its sharp-frequency
expression and the reference interval average in the definition
of $\Xi$. The first error, multiplied by $m^2/(2h_B)$, is at most
$D_{\rm all}^{1}m^2/(\pi h_BS_{m,B})$ by
Corollary~\ref{cor:ui-normalization}. The difference between the
normalized actual mass and the reference at its center is bounded
by $\mathscr E_m(h_B)$. The difference between that reference
value and its interval average is at most $h_BD_m/2$, by
Lemma~\ref{lem:ui-reference-modulus}. This proves
\eqref{eq:ui-xi-inherited}. The stated bandwidth makes its last
term $D_{\rm all}^{1}m^{-P}/\pi$. Every term therefore vanishes
with $B$ fixed.

For the sharper raw implication, no frequency truncation is
needed. At each Lebesgue point of the exact-label BV density,
insert its exact interval average directly between the density
and the reference at the center. The local-variation estimate in
the proof of Proposition~\ref{prop:ui-raw-bridge} gives
\[
 \sup_{R,n,k}\esssup_{t\in\mathcal C_{m,R,n,k}(M)}
  |m^2p_{n,R}(k,m,t)-\mathcal L_{m,R}(k_1,k_2,t,n)|
       \le\Omega^p_{m,B}+\mathscr E_m(h_B).
\]
The inherited fixed-window theorem is uniform over all the
displayed radii and labels. Take the collision limsup, keeping
$h_B$ fixed, to obtain \eqref{eq:ui-local-charge-suffices}; only
then take the outer limit. Exceptional Lebesgue-point sets may
depend on the radius, as allowed by the original
$\sup_R\esssup_t$ order.
\end{proof}

This argument does not claim a new strong-operator estimate on a
growing roof band. It evaluates the particular scalar sharp
interval expression through the retained fixed-window law and
the proved tail. A direct growing-band spectral estimate is
another possible method, not a prerequisite for this implication.
The analysis intervals do not replace the original microscopic
conditioning events. The conclusion in this subsection is for
$w=1,V=1$; its use for additional observation classes requires
their own stated fixed-window theorem.

\paragraph{Why the local variation condition cannot be omitted.}
The following analytical example tests the implication only; it is
not substituted for a billiard source.  Fix $\tau_0\in(\tau_-,\tau_+)$
and, for large $m$, put
\[
 q_m(t)=m^{-2}\mathbf1_{[m\tau_0,m\tau_0+e^{-m}]}(t),\qquad g_m=0.
\]
It is a nonnegative subprobability with the correct enclosing
physical support, $\|Dq_m\|_{\rm TV}=2m^{-2}$, and
$\|q_m\|_1=m^{-2}e^{-m}$.  Every interval of a fixed positive
length $2h$ satisfies
$m^2(2h)^{-1}\int_Iq_m\le e^{-m}/(2h)$, uniformly in its location.
Nevertheless $m^2\|q_m\|_\infty=1$.  Moreover
$|\widehat q_m(s)|\le m^{-2}e^{-m}$, so its sharp-frequency
window integral satisfies
\[
 {m^2\over2h}|\mathfrak I_S(I)|
       \le {S e^{-m}\over\pi}.
\]
It thus has vanishing normalized finite-band error at every
polynomial bandwidth, as well as the common BV bound and vanishing
normalized interval masses.  Its local derivative charge, however,
remains of order $m^{-2}$ at the spike.  This verifies that the new
absolute inversion and a low- or polynomial-frequency comparison
cannot alone be promoted to an essential-height assertion.  A
physical mechanism paying the actual local charge or the original
same-roof shortage is still needed.  The written fixed-band joint
spectral, arithmetic and resonance arguments in historical native
84--87 are retained at their stated scope.  The present bandwidth
grows with the collision count and the reciprocal microscopic width;
no new strong-operator estimate on that growing band has been proved.
Nevertheless Proposition~\ref{prop:ui-inherited-bridge} evaluates
the scalar analysis-window expression through the retained
fixed-window law and leaves the complete normalized local charge
as a sufficient closing condition on that route. The exact
arithmetic reference, other observation classes, microscopic
conditional denominators and long-time local inverse retain their
own input-specific obligations.
